% capitolo 1: introduzione
\chapter{Contesto aziendale}
\label{cap:contesto-aziendale}

\intro{Il presente capitolo ha lo scopo di introdurre il contesto aziendale in cui si è svolta l'esperienza di stage. Verrà presentata la realtà di Wintech S.p.A. esplorandone 
l'offerta, l'approccio alla sicurezza informatica e il modello organizzativo.}\\

\section{L'azienda}
Wintech (Winning Technologies S.p.A.) è una consolidata realtà italiana operante nel settore dell'Information Technology.\footcite{wintech:website}
Fondata nel 1987, l'azienda ha la sua sede principale a Padova e presidia strategicamente il territorio con ulteriori filiali operative a Milano e Bassano del Grappa.

Ad oggi l'azienda conta un organico di oltre 95 professionisti, un portafoglio di più di 500 clienti attivi e un fatturato superiore ai 12 milioni di euro.
Wintech si posiziona sul mercato principalmente come consulente e \emph{System Integrator}: il suo core business consiste nel combinare competenze e tecnologie avanzate per far comunicare in 
sinergia molteplici sottosistemi eterogenei, creando un'infrastruttura IT coesa, sicura e funzionale che porti un reale valore aggiunto al business del cliente.

I prodotti e i servizi offerti si rivolgono a un bacino d'utenza ampio e diversificato, studiato per rispondere alle esigenze di liberi professionisti, Piccole e Medie Imprese (PMI), 
grandi aziende, Pubblica Amministrazione, istituti bancari e assicurativi. 

Per garantire lo sviluppo e l'integrazione di soluzioni all'avanguardia, stabili e scalabili, Wintech ha stretto nel tempo solide \emph{partnership} con i principali colossi mondiali 
dell'informatica, tra cui Microsoft, IBM, Oracle, HP e Sophos. Tali collaborazioni strategiche garantiscono l'accesso diretto alle migliori tecnologie sul mercato e certificano 
l'elevato livello di competenza raggiunto dal personale aziendale.

Infine, un aspetto cardine dell'identità aziendale è la profonda attenzione rivolta alla sicurezza dei dati e alla qualità dei processi interni. 
A testimonianza di questo impegno, Wintech è ufficialmente dotata delle prestigiose certificazioni internazionali ISO 9001, inerente ai sistemi di gestione per la qualità, 
e ISO/IEC 27001, riguardante i sistemi di gestione per la sicurezza delle informazioni.


% logo Wintech



\section{L'offerta}
L'offerta di Wintech si articola in due macro-aree complementari: i servizi di consulenza e progettazione IT, e la fornitura di applicazioni gestionali specifiche per i diversi settori di mercato.

\subsection{Servizi}
L'azienda accompagna i clienti in ogni fase del loro percorso di transizione digitale. I servizi offerti spaziano dall'analisi e progettazione di soluzioni su misura, 
al \emph{deployment} e all'avviamento dei sistemi, fino alla formazione del personale. Viene inoltre garantita una costante assistenza tecnica e applicativa, sia tramite \emph{help desk} 
remoto che con interventi diretti sul posto. Un pilastro fondamentale dei servizi Wintech è la consulenza in ambito Cyber Security e architetture Cloud, per garantire la protezione, 
il monitoraggio e il Disaster Recovery delle infrastrutture di rete.

\subsection{Prodotti e applicazioni gestionali}
Wintech personalizza e fornisce un ampio portafoglio di \emph{software}, adattandoli alle dimensioni e alle esigenze operative del cliente. L'offerta si suddivide principalmente in:
\begin{itemize}
    \item \textbf{Professionisti e studi}: Soluzioni integrate per l'amministrazione, la fatturazione, la contabilità e la gestione avanzata del personale (tra cui le suite PROFIS, JOB e STUDIO).
    \item \textbf{PMI e grandi aziende}: Sistemi avanzati di Enterprise Resource Planning (ERP) per il controllo di gestione (come eSOLVER), strumenti di Business Intelligence per l'analisi dei dati, e piattaforme per il Customer Relationship Management (CRM).
    \item \textbf{Aziende di moda}: Declinazioni specifiche dei gestionali (es. eSOLVER per la Moda) sviluppate per gestire le varianti di prodotto, la logistica, le spedizioni e i punti vendita tipici del settore.
    \item \textbf{Banche e assicurazioni}: Soluzioni dedicate alla sicurezza documentale, all'archiviazione e all'integrazione della firma grafometrica.
\end{itemize}

A completamento dell'offerta, l'azienda progetta e sviluppa portali web e piattaforme di \emph{Collaboration} aziendale (intranet) per favorire la comunicazione interna e la condivisione delle informazioni.

\subsection{Partner tecnologici}
Per garantire standard di qualità eccellenti e fornire tecnologie sempre all'avanguardia, Wintech si avvale della collaborazione dei principali colossi mondiali del settore IT. 
Le solide \emph{partnership} instaurate nel tempo permettono all'azienda di accedere in modo diretto a soluzioni certificate, scalabili e sicure.

\begin{figure}[H]
    \centering
    % loghi dei partner principali
    \includegraphics[width=1\textwidth]{partner.png}
    \vspace{0.2cm}
    \caption{Principali \emph{partner} di Wintech}
    \label{fig:aziende_partner}
\end{figure}

\section{Certificazioni aziendali}
Wintech è un'azienda certificata ISO 27001 e ISO 9001, che garantiscono massimi standard di qualità professionale e di gestione della sicurezza ai collaboratori, nonchè reali vantaggi 
competitivi ai clienti.
Tutte le certificazioni:
ISO 9001
ISO IEC 27001
ISO IEC 27017
ISO IEC 27018
ISO 22301
% da descrivere in molto breve

% ISO 9001 (Qualità) e ISO/IEC 27001 (Sicurezza delle informazioni)
% immagine/infografica sulla sicurezza dei dati (da prendere sul sito)


\section{Organizzazione e processi interni}
% organizzazione aziendale: team ()

% processi e way of working interno
% metodologia agile (framework scrum) - con immagine, microsoft teams, regole sulla comunicazione
% Un aspetto positivo dello stage è stata la comunicazione con i colleghi e con la tutor interna, che mi ha fornito supporto quotidiano e con cui ci sono stati confronti aperti



\section{Organizzazione del testo}
Il presente documento è strutturato in nove capitoli, organizzati come segue:

\begin{itemize}
    \item \textbf{Il primo capitolo} introduce il contesto aziendale di Wintech S.p.A..
    \item \textbf{Il secondo capitolo} fornisce una descrizione dettagliata dell'esperienza di stage.
    \item \textbf{Il terzo capitolo} si focalizza sullo studio dei casi d'uso e l'analisi dei requisiti aziendali.
    \item \textbf{Il quarto capitolo} analizza e confronta la piattaforma OpenMetadata con i principali competitor sul mercato in ambito Data Catalog.
    \item \textbf{Il quinto capitolo} è dedicato all'implementazione tecnica, illustrando le fasi di installazione, configurazione e Data Discovery.
    \item \textbf{Il sesto capitolo} approfondisce i concetti di Data Quality e Data Observability per il monitoraggio dello stato di salute dei dati.
    \item \textbf{Il settimo capitolo} tratta il tema della Data Governance e dello sviluppo sicuro.
    \item \textbf{L'ottavo capitolo} approfondisce l'integrazione tra il Data Catalog e gli strumenti di Business Intelligence, in particolare Microsoft Power BI.
    \item \textbf{Il nono capitolo} presenta le conclusioni, con un consuntivo finale sul raggiungimento degli obiettivi, le problematiche affrontate e le competenze acquisite.
\end{itemize}

\subsection{Convenzioni}
Riguardo la stesura del testo, relativamente al documento sono state adottate le seguenti convenzioni tipografiche:
\begin{itemize}
	\item gli acronimi, le abbreviazioni e i termini ambigui o di uso non comune menzionati vengono definiti nel glossario, situato alla fine del presente documento.
	\item per la prima occorrenza dei termini riportati nel glossario viene utilizzata la seguente nomenclatura: \emph{parola}\glsfirstoccur.
	\item particolari termini in lingua straniera o facenti parti del gergo tecnico sono evidenziati con il carattere \emph{corsivo}. Fanno eccezione i macro-concetti o i domini disciplinari del progetto (come Data Quality, Data Catalog, Data Lineage...), 
    i quali vengono scritti in tondo con le iniziali maiuscole per agevolare la scorrevolezza della lettura.
\end{itemize}
